\section{Discrete Random Variables}

We in this section study the probability spaces of the form \(\para{\Omega, \mathcal{F}, \Prob}\) where
\[
    \Omega = \brce{\omega_1, \omega_2, \ldots} \quad \text{and} \quad \mathcal{F} = \powerset\para{\Omega}.
\]

The reason we study this countable case is as follows. If we know the probability of all the singletons \(\Prob\para{\brce{\omega_i}}\), then for all \(A \subseteq \Omega\), we have
\[
    \Prob\para{A} = \Prob\para{\bigcup_{\omega_i \in A} \brce{\omega_i}} = \sum_{\omega_i \in A} \Prob\para{\brce{\omega_i}}.
\]

This allows us to pinpoint the distribution just by investigating the singletons.

\subsection{Discrete Probability Distribution}

\begin{definition}[Discrete Probability Distribution]
    \label{def:discrete-probability-distribution}%
    We call \(\para{\Prob\para{\brce{\omega_i}}}_i\) a \emph{discrete probability distribution}, and we write
    \[
        p_i = \Prob\para{\brce{\omega_i}}
    \]
    for all \(i\).
\end{definition}

We have:
\begin{itemize}
    \item \(p_i \geq 0\) for all \(i\), and
    \item \(\sum_i p_i = 1\).
\end{itemize}

\subsection{Common Discrete Probability Distributions}

Now we see some common examples of discrete probability distributions.

\begin{example}[Tossing a Biased Coin]
    \label{ex:tossing-biased-coin}%
    We toss a biased coin with probability \(p\) landing on its head and probability \(\para{1 - p}\) landing on its tail. Then we may model this as
    \[
        \Prob\para{\text{landing on tail}} = 1 - p \quad \text{and} \quad \Prob\para{\text{landing on head}} = p.
    \]
\end{example}

This is modelled by the simplest non-trivial discrete probability distribution, the Bernoulli distribution.

\begin{definition}[Bernoulli Distribution]
    \label{def:bernoulli-distribution}%
    The \emph{Bernoulli distribution}, denoted as \(\Bernoulli\para{p}\) for \(p \in \ccintv{0}{1}\), is defined as follows. The sample space is \(\Omega = \brce{0, 1}\), and the distribution is defined by
    \[
        p_0 = \Prob\para{\brce{0}} = 1 - p \quad \text{and} \quad p_1 = \Prob\para{\brce{1}} = p.
    \]
\end{definition}

Now, what happens if we toss multiple \(p\)-coins? 

\begin{example}[Tossing Multiple \(p\)-coins]
    \label{ex:tossing-multiple-biased-coins}%
    Now, we toss \(n\) such biased coins. Then we may model this as
    \[
        \Prob\para{\text{landing on heads }k\text{ times}} = \binom{n}{k} p^k \para{1 - p}^{n - k}.
    \]
\end{example}

\begin{definition}[Binomial Distribution]
    \label{def:binomial-distribution}%
    The \emph{Binomial distribution}, denoted as \(\Binomial\para{n, p}\), for \(n \in \NN\) and \(p \in \ccintv{0}{1}\), is defined as follows. The sample space is \(\Omega = \brce{0, 1, \ldots, n}\), and the distribution is defined by
    \[
        p_k = \Prob\para{\brce{k}}= \binom{n}{k} p^k \para{1 - p}^{n - k}
    \]
    where \(k = 0, 1, \ldots, n\).
\end{definition}

This could be thought of the sum of \emph{independent} Bernoulli random variables as well, as we will see later. Indeed, when \(n = 1\), the Binomial distribution is simply the Bernoulli distribution.

The Binomial distribution is a generalisation of the Bernoulli distribution by having multiple trials. Similarly, we can generalise the Binomial distribution by introducing multiple possible outcomes.

\begin{example}[Throwing Balls into Boxes]
    \label{ex:throwing-balls-into-boxes}%
    We throw \(n\) balls into \(k\) boxes, landing into each box with probability \(p_1, p_2, \ldots, p_k\) with \(\sum p_i = 1\) and \(p_i \in \ccintv{0}{1}\). Then, the probability of having \(n_i\) balls in each of the \(k\) boxes, where \(\sum n_i = n\), is given by
    \[
        \binom{n}{n_1, n_2, \ldots, n_k} p_1^{n_1} p_2^{n_2} \cdots p_k^{n_k}
    \]
    using the multinomial coefficient.
\end{example}

Alternatively, we can also model this as a polyhedron landing on different faces, with unequal probability (as opposing to a coin only having two sides). Although, obviously, there are no polyhedrons with at most 3 faces, but this is probably good enough.

\begin{example}[Polyhedron landing on Different sides]
    \label{ex:polyhedron-different-sides}%
    We have a \(k\)-sided polyhedron with a probability of \(p_i\) landing on side \(i\). Throwing the polyhedron \(n\) times, the probability that it being landing on side \(i\) for \(n_i\) times respectively, for \(\sum n_i = n\), gives the same probability distribution as above.
\end{example}

\begin{definition}[Multinomial Distribution]
    \label{def:multinomial-distribution}%
    The \emph{Multinomial distribution}, denoted as \(\Multinomial\para{n, p_1, \ldots, p_k}\), with \(p_i \geq 0\) and \(\sum_{i = 1}^{k} p_i = 1\), is defined as follows. The sample space is
    \[
        \Omega = \set{\para{n_1, n_2, \ldots, n_k} \in \NN^k}{\sum_{i = 1}^{k} n_i = n},
    \]
    and the distribution is given by
    \[
        p_{\para{n_1, n_2, \ldots, n_k}} = \binom{n}{n_1, n_2, \ldots, n_k} p_1^{n_1} p_2^{n_2} \cdots p_k^{n_k},
    \]
    for \(\sum_{i = 1}^{k} n_i = n\).
\end{definition}

So far, our sample space has been finite. But we can also define discrete distributions over infinite sample spaces.

\begin{example}[Tossing Coins until a Head]
    \label{ex:tossing-coins-until-head}%
    We toss a \(p\)-coin, having result Head with probability \(p\). We toss it independently, until we see the first Head. The probability of the event of tossing \(k\) coins until the first head (inclusive), is given by
    \[
        \Prob\para{\text{the first head is the }k\text{th coin}} = \para{1 - p}^{k - 1} \cdot p.
    \]
\end{example}

\begin{definition}[Geometric Distribution]
    \label{def:geometric-distribution}%
    The \emph{Geometric distribution}, denoted as \(\Geometric\para{p}\), with \(p \in \ccintv{0}{1}\), is defined as follows. The sample space is
    \[
        \Omega = \brce{1, 2, \ldots},
    \]
    and the distribution is given by
    \[
        p_k = \para{1 - p}^{k - 1} \cdot p,
    \]
    for \(k \in \Omega\).
\end{definition}

There is a slight caveat to the geometric distribution -- some people prefer modelling the distribution starting from \(0\), i.e., the number of toss \emph{before, but not including} the final toss getting the head. In this case, the sample space would start from zero, given by
\[
    \Omega' = \brce{0, 1, 2, \ldots}
\]
and the probability distribution is given by
\[
    p'_k = \para{1 - p}^{k} \cdot p.
\]

It would be nice to specify whether a geometric distribution starts from \(0\) or from \(1\).

The final distribution to introduce is not very intuitive when first seeing the distribution expression -- but the following example gives a motivation.

\begin{example}[Constant Average Rate]
    We suppose that customers arrive into a shop in \(\ccintv{0}{1}\) at a constant rate, with an average of \(\lambda > 0\) customers arriving. To model this, we discretise \(\ccintv{0}{1}\) into a collection of sub-intervals
    \[
        \set{\ccintv{\frac{i - 1}{N}}{\frac{i}{N}}}{i = 1, 2, \ldots, N},
    \]
    and in each interval we toss a \(p\)-coin to decide whether a customer would arrive, independently. In this sub-interval approximation, the probability that \(k\) customers have arrived is given by the binomial distribution \(\Binomial\para{N, p}\), so having a probability of
    \[
        \binom{N}{k} \cdot p^k \cdot \para{1 - p}^{N - k}
    \]
    for \(k = 0, 1, 2, \ldots, N\).

    By our model, assuming there is an average of \(\lambda\) customers arriving in the whole interval, so in each sub-interval, we take \(p = \frac{\lambda}{N}\). (We may take \(N\) to be big such that \(p < 1\).) The shorter the intervals split up, the better this models the continuous arrivals. So we fix \(k\) and let \(N \to \infty\), giving the probabilities
    \[
        \Prob\para{k\text{ customers arrive}} = \lim_{N \to \infty} \frac{N!}{k! \para{N - k}!} \cdot \para{\frac{\lambda}{N}}^k \cdot \para{1 - \frac{\lambda}{N}}^{N - k}.
    \]

    Evaluating this limit gives us
    \begin{align*}
        p_k &= \lim_{N \to \infty} \brac{\frac{\lambda^k}{k!} \cdot \frac{N \para{N - 1} \cdots \para{N - k + 1}}{N^k} \para{1 - \frac{\lambda}{N}}^{N - k}}\\
        &= \frac{\lambda^k}{k!} \cdot 1 \cdot e^{-\lambda} \\
        &= e^{-\lambda} \cdot \frac{\lambda^{k}}{k!},
    \end{align*}
    and this is indeed a probability distribution since
    \[
        \sum_{k = 0}^{\infty} p_k = 1
    \]
\end{example}

\begin{definition}[Poisson Distribution]
    \label{def:poisson-distribution}%
    The \emph{Poisson distribution}, denoted as \(\Poisson\para{\lambda}\), for parameter \(\lambda > 0\), is defined as follows. The sample space is
    \[
        \Omega = \brce{0, 1, 2, \ldots},
    \]
    and the distribution is given by
    \[
        p_k = e^{-\lambda} \cdot \frac{\lambda^k}{k!}
    \]
    for \(k \in \Omega\).
\end{definition}

The Poisson distribution can be used to model the number of occurrences of an event in a given interval of time.

\subsection{Random Variable}

\begin{definition}[Random Variable]
    \label{def:random-variable}%
    For a probability space \(\para{\Omega, \mathcal{F}, \Prob}\), a \emph{random variable} is a function \(\functiondomain{X}{\Omega}{\RR}\) with the property that for all \(x \in \RR\), we have
    \[
        \set{\omega \in \Omega}{X\para{\omega} \leq x} = X\inverse \para{\ocintv{-\infty}{\omega}} \in \mathcal{F}.
    \]
\end{definition}

\begin{notation}
    As an abuse of notation, we denote for all \(x \in \RR\),
    \[
        \brce{X \leq x} = \set{\omega \in \Omega}{X\para{\omega} \leq x}
    \]
    and for all \(A \subseteq \RR\)
    \[
        \brce{X \in A} = X\inverse\para{A} = \set{\omega \in \Omega}{X\para{\omega} \in A}.
    \]
\end{notation}

\begin{definition}[Indicator Random Variable]
    \label{def:indicator-random-variable}%
    The \emph{indicator} of a set \(A \in \mathcal{F}\) is denoted as
    \[
        \indicator_{A} \para{\omega} = \indicator\para{\omega \in A} = \begin{cases}
            1, & \omega \in A, \\ 0, & \omega \in A\compl.
        \end{cases}
    \]

    Since \(A \in \mathcal{F}\), \(\indicator_{A}\) is indeed a random variable, and we may restrict the co-domain to
    \[
        \functiondomain{\indicator_{A}}{\Omega}{\brce{0, 1}}.
    \]
\end{definition}


\begin{definition}[Probability Distribution Function]
    \label{def:probability-distribution-function}%
    The \emph{probability distribution function} or \emph{cumulative distribution function} of a random variable \(X\) is, for \(x \in \RR\),
    \[
        \function{F_X}{\RR}{\ccintv{0}{1}}{x}{\Prob\para{X \leq x}.}
    \]
\end{definition}

\begin{definition}[Random Variable in \(\RR^n\)]
    \label{def:random-variable-rn}%
    The combined function
    \[
        X = \para{X_1, X_2, \ldots, X_n}
    \]
    is called a \emph{random variable in \(\RR^n\)}, if
    \[
        \functiondomain{\para{X_1, X_2, \ldots, X_n}}{\Omega}{\RR^n},
    \]
    with the additional restriction that for all \(x_1, x_2, \ldots, x_n \in \RR\), we have
    \begin{align*}
        &\phantom{=} X\inverse\para{\ocintv{-\infty}{x_1} \times \ocintv{-\infty}{x_2} \times \cdots \times \ocintv{-\infty}{x_n}} \\
        &= \set{\omega \in \Omega}{X_1 \para{\omega} \leq x_1 \land X_2 \para{\omega} \leq x_2 \land \cdots \land X_n \para{\omega} \leq x_n} \\
        &\in \mathcal{F}.
    \end{align*}

    Equivalently, this restriction can be stated as that all \(X_i\)s are real random variables, since
    \[
        \brce{X_1 \leq x_1, \ldots, X_n \leq x_n} = \brce{X_1 \leq x_1} \cap \brce{X_2 \leq x_2} \cap \cdots \cap \brce{X_n \leq x_n}.
    \]
\end{definition}

We now restrict back our view to discrete probability distributions.

\begin{definition}[Discrete Random Variable]
    \label{def:discrete-random-variable}%
    A random variable \(X\) is called \emph{discrete} if it takes values in a countable set, i.e., \(\img \para{X} \subseteq \RR\) is countable.
\end{definition}

\begin{definition}[Probability Mass Function]
    \label{def:probability-mass-function}%
    Suppose a discrete random variable \(X\) satisfies \(\img X = S \subseteq \RR\), which is countable. For every \(x \in S\), we denote
    \[
        p_x = \Prob\para{X = x} = \Prob\para{\set{\omega \in \Omega}{X\para{\omega = x}}},
    \]
    and we call \(\para{p_x}_{x \in S}\) the \emph{probability mass function} or the \emph{distribution} of \(X\).
\end{definition}

From the definition of a probability mass function, we have for all \(A \subseteq S\), that
\[
    \Prob\para{X \in A} = \sum_{x \in A} p_x.
\]

If the \(p_x\)s are the Bernoulli distribution, then we say \(X\) is a Bernoulli \emph{random variable}, and it has the Bernoulli \emph{distribution}. Similar terminology applies for the other distributions as well.

Recall that two events \(A, B\) are independent, if
\[
    \Prob \para{A \cap B} = \Prob\para{A} \cdot \Prob\para{B}.
\]

We may extend this notion to random variables as well.

\begin{definition}[Independence of Random Variables]
    \label{def:independence-random-variable}%
    Let \(X_1, X_2, \ldots, X_n\) be discrete random variables, satisfying \(\img X_i = S_i\) for \(i = 1, 2, \ldots, n\). Then they are \emph{independent}, if for any \(x_i \in S_i\) for \(i = 1, 2, \ldots, n\), that
    \[
        \Prob\para{X_1 = x_1, X_2 = x_2, \ldots, X_n = x_n} = \Prob\para{X_1 = x_1} \cdot \Prob\para{X_2 = x_2} \cdots \Prob\para{X_n = x_n},
    \]
    where the left-hand side notation denotes the intersection of all the events.
\end{definition}

In particular, for two random variables \(X_1, X_2\), they are independent if
\[
    \Prob\para{X_1 = x_1, X_2 = x_2} = \Prob\para{X_1 = x_1} \cdot \Prob\para{X_2 = x_2},
\]
for all \(x_1 \in \img X_1, x_2 \in \img X_2\).

\begin{example}[Independence of Bernoulli Random Variables]
    \label{ex:independence-bernoulli}
    We toss a \(p\)-coin \(N\) independently, with
    \[
        \Omega = \brce{0, 1}^N.
    \]

    Let \(\omega = \para{\omega_1, \ldots, \omega_N} \in \Omega\), where \(\omega_i\) be the outcome of the \(i\)-th toss. We hence have
    \[
        p_\omega = \Prob\para{\brce{\omega}} = \prod_{k = 1}^{N} p^{\omega_k} \para{1 - p}^{1 - \omega_k}
    \]

    For all \(k = 1, 2, \ldots, N\), we define \(\functiondomain{X_k}{\Omega}{\brce{0, 1}}\), with \(X_k\para{\omega} = \omega_k\). For example, we have
    \begin{align*}
        \Prob\para{X_k = 1} &= p = \Prob\para{\omega_k = 1} \\
        \Prob\para{X_k = 0} &= 1 - p = \Prob\para{\omega_k = 0}.
    \end{align*}

    We have \(X_k \sim \Bernoulli\para{p}\). Indeed, \(X_1, X_2, \ldots, X_N\) are independent random variables, since for all \(x_1, x_2, \ldots, x_N \in \brce{0, 1}\),
    \begin{align*}
        \Prob\para{X_1 = x_1, X_2 = x_2, \ldots, X_N = x_N} &= \prod_{k = 1}^{N} p^{x_k} \para{1 - p}^{1 - x_k} \\
        &= \prod_{i = 1}^{N} \Prob\para{X_i = x_i}.
    \end{align*}

    We let the sum of the random variables be
    \[
        S_N\para{\omega} = X_1 \para{\omega} + \cdots + X_N \para{\omega} = \text{number of heads in }N\text{ tosses}.
    \]

    Therefore,
    \[
        \Prob\para{S_N = k} = \binom{N}{k} \cdot p^k \cdot \para{1 - p}^{N - k}
    \]
    and we therefore see \(S_N \sim \Binomial\para{N, p}\) -- sum of independent Bernoulli random variables is a Binomial random variable.
\end{example}

\subsection{Expectation}

Let \(\Omega\) be finite or uncountable, and let \(\functiondomain{X}{\Omega}{\RR}\) be a discrete random variable.

We say \(X\) is \emph{non-negative} if \(X \geq 0\). We first define the expectation for a non-negative random variable (see later for motivation).

\begin{definition}[Expectation of Non-Negative Random Variable]
    \label{def:expectation-non-negative}%
    The \emph{expectation}, or \emph{mean/average} of a non-negative random variable is given by
    \[
        \Expt \brac{X} = \sum_{\omega \in \Omega} \Prob\para{\brce{\omega}} \cdot X\para{\omega}.
    \]
\end{definition}

Equivalently, we may define the expectation as follows. If we let the range of the random variable be \(\Omega_X = \set{X\para{x}}{\omega \in \Omega}\), and partition the sample space as
\[
    \Omega = \bigcup_{x \in \Omega_X} \brce{X = x},
\]
then the expectation satisfies
\begin{align*}
    \Expt\brac{X} &= \sum_{\omega \in \Omega} X\para{\omega} \Prob\para{\brce{\omega}} \\
    &= \sum_{x \in \Omega_X} \sum_{\omega \in \brce{X = x}} X\para{\omega} \cdot \Prob\para{\brce{\omega}} \\
    &= \sum_{x \in \Omega_X} \sum_{\omega \in \brce{X = x}} x \cdot \Prob\para{\brce{\omega}} \\
    &= \sum_{x \in \Omega_X} x \sum_{\omega \in \brce{X = x}} \cdot \Prob\para{\brce{\omega}} \\
    &= \sum_{x \in \Omega_X} x \cdot \Prob\para{X = x}.
\end{align*}

In other words, the expectation of \(X\) is the average of the values \(x\) taken by \(X\), with weights given by the corresponding probabilities \(\Prob\para{X = x}\).

\begin{example}[Expectation of Binomial Random Variable]
    \label{ex:expectation-binomial}%
    Let \(X \sim \Binomial\para{N, p}\). This has probability mass function
    \[
        \Prob\para{X = k} = \binom{N}{k} p^k \para{1 - p}^{N - k}.
    \]

    By definition,
    \begin{align*}
        \Expt\brac{X} &= \sum_{k = 0}^{N} k \Prob\para{X = k} \\
        &= \sum_{k = 0}^{N} k \binom{N}{k} p^k \para{1 - p}^{N - k} \\
        &= \sum_{k = 0}^{N} k \cdot \frac{N!}{\para{N - k}! k!} p^k \para{1 - p}^{N - k} \\
        &= \sum_{k = 1}^{N} N \cdot \frac{\para{N - 1}!}{\para{N - k}! \para{k - 1}!} p^k \para{1 - p}^{N - k} \\
        &= \sum_{k = 0}^{N - 1} N \binom{N - 1}{k} p^{k} \para{1 - p}^{N - k - 1} p \\
        &= Np \para{p + 1 - p}^{N - 1} \\
        &= Np.
    \end{align*}
\end{example}

\begin{example}[Expectation of Poisson Random Variable]
    \label{ex:expectation-poisson}%
    Let \(X \sim \Poisson\para{\lambda}\) with \(\lambda > 0\). We have
    \[
        \Prob\para{X = k} = e^{-\lambda} \cdot \frac{\lambda^k}{k!}.
    \]

    By definition,
    \begin{align*}
        \Expt\brac{X} &= \sum_{k = 0}^{\infty} k \cdot \Prob\para{X = k} \\
        &= \sum_{k = 1}^{\infty} k \cdot e^{-\lambda} \cdot \frac{\lambda^k}{k!} \\
        &= e^{-\lambda} \cdot \lambda \cdot \sum_{k = 1}^{\infty} \frac{\lambda^{k - 1}}{\para{k - 1}!} \\
        &=\lambda.
    \end{align*}
\end{example}

Now we define the expectation of a general random variable \(X\). The reason why we defined it for a non-negative random variable first, is due the sum being over a \emph{set}, which has no built-in ordering to specify the sum. It is always fine to rearrange a non-negative sum, since it either absolutely converges or always diverges, and rearranging does not change its sum.

However, if we straight-away define it for a general random variable using the same sum, and if the series is only \emph{conditionally convergent}, then changing the order of summation will change the limit, which makes the expectation change. Therefore, we will formalise this considering the sum of the positive and negative values separately.

Let \(X^+ = \max\brce{X, 0}\) and \(X^- = \min\brce{-X, 0}\).

We therefore have \(X = X^+ - X^-\) and \(\abs{X} = X^+ + X^-\).

Suppose \(X\) is discrete. Using \autoref{def:expectation-non-negative}, we have already defined \(\Expt\brac{X^+}\) and \(\Expt\brac{X^-}\).

\begin{definition}[Expectation of General Random Variable]
    \label{def:expectation-general-random-variable}%
    Let \(X\) be a discrete random variable and \(X^+\) and \(X^-\) be defined as above. If at least one of \(\Expt\brac{X^+}\) and \(\Expt\brac{X^-}\) is finite, then we define the \emph{expectation} of \(X\) as
    \[
        \Expt\brac{X} = \Expt\brac{X^+} - \Expt\brac{X^-}.
    \]

    If both of them are infinite, then the expectation of \(X\) is \emph{undefined}.
\end{definition}

When we write \(\Expt\brac{X}\), we implicitly imply that it is defined.

\begin{definition}[Integrablility of Random Variable]
    \label{def:integrability-random-variable}%
    If \(\Expt\brac{\abs{X}} < \infty\), then we say \(X\) is \emph{integrable}.
\end{definition}

When \(\Expt\brac{X}\) is well-defined, then we have
\[
    \Expt\brac{X} = \sum_{x \in \Omega_X} x \Prob\para{X = x}.
\]

\begin{proposition}[Properties of Expectation]
    \label{prop:properties-expectation}%
    \begin{enumerate}
        \item If \(X \geq 0\), then \(\Expt\brac{X} \geq 0\).
        \item If \(X \geq 0\) and \(\Expt\brac{X} = 0 \), then \(\Prob\para{X = 0} = 1\).
        \item If \(c \in \RR\), then \(\Expt\brac{cX} = c \Expt\brac{X}\), and \(\Expt\brac{c + X} = c + \Expt\brac{X}\).
        \item If \(X\) and \(Y\) are random variables, then \(\Expt\brac{X + Y} = \Expt{X} + \Expt{Y}\),
        \item Let \(c_1, \ldots, c_n \in \RR\), and \(X_1, \ldots, X_n\) be integrable random variables. Then we have
        \[
            \Expt \brac{\sum_{i = 1}^{n} c_i X_i} = \sum_{i = 1}^{n} c_i \Expt \brac{X_i}.
        \]
        \item Suppose \(X_1, X_2, \ldots\) are non-negative random variables. Then
        \[
            \Expt\brac{\sum_{n \in \NN} X_n} = \sum_{n \in \NN} \Expt\brac{X_n}.
        \]
    \end{enumerate}
\end{proposition}

\begin{proof}[on Linearity over Countable Collection]
    We have
    \begin{align*}
        \Expt\brac{\sum_{n \in \NN} X_n} &= \sum_{\omega \in \Omega} \para{\sum_{n \in \NN} X_n \para{\omega}} \cdot \Prob\para{\brce{\omega}} \\
        &= \sum_{n \in \NN} \sum_{\omega \ni \Omega} X_n \para{\omega} \Prob\para{\brce{\omega}} \\
        &= \sum_{n \in \NN} \Expt\brac{X_n}.
    \end{align*}
\end{proof}

\begin{example}[Expectation of Indicator]
    \label{ex:expectation-indicator}%
    Recall the definition of the \emph{indicator} of an event \(A \in \mathcal{F}\) is \(\indicator\para{A}\). Let \(X = \indicator\para{A}\), then \(X\para{\omega} = \indicator\para{\omega \in A}\). Therefore, \(\Expt\brac{X} = \Prob\para{A}\).
\end{example}

\begin{proposition}[Expectation of Function of Random Variable]
    \label{prop:expectation-function-random-variable}%
    Let \(\functiondomain{X}{\Omega}{\RR}\) be a random variable, and let \(\functiondomain{g}{\RR}{\RR}\). We define
    \[
        g\para{X} \para{\omega} = g\para{X\para{\omega}}.
    \]

    Then
    \[
        \Expt\brac{g\para{X}} = \sum_{x \in \Omega_X} g\para{x} \cdot \Prob\para{X = x}.
    \]
\end{proposition}

\begin{proof}
    Set \(Y = g\para{X}\). Then we have
    \[
        \Expt\brac{Y} = \sum_{y \in \Omega_Y} y \cdot \Prob\para{Y = y}.
    \]

    Now, note that the event can be simplified as
    \begin{align*}
        \brce{Y = y} &= \set{\omega \in \Omega}{Y \para{\omega} = y} \\
        &= \set{\omega \in \Omega}{g\para{X\para{\omega}} = y} \\
        &= \set{\omega \in \Omega}{X\para{\omega} \in g\inverse \para{\brce{y}}} \\
        &= \brce{X \in g\inverse \para{\brce{y}}}.
    \end{align*}

    Therefore,
    \begin{align*}
        \Expt\brac{Y} &= \sum_{y \in \Omega_Y} y \cdot \Prob\para{X \in g\inverse\para{\brce{y}}} \\
        &= \sum_{y \in \Omega_Y} y \cdot \sum_{x \in g\inverse\para{\brce{y}}} \Prob\para{X = x} \\
        &= \sum_{y \in \Omega_Y} \sum_{x \in g\inverse\para{\brce{y}}} g\para{x} \cdot \Prob\para{X = x} \\
        &= \sum_{x \in \Omega_X} g\para{x} \Prob\para{X = x}.
    \end{align*}
\end{proof}

\begin{proposition}[Expectation of Natural Number-Values Random Variables]
    \label{prop:expt-natural-number-random-variable}%
    Suppose \(X \geq 0\) and takes integer values. Then we have
    \[
        \Expt\brac{X} = \sum_{k = 1}^{\infty} \Prob\para{X \geq k} = \sum_{k = 0}^{\infty} \Prob\para{X > k}.
    \]
\end{proposition}

\begin{proof}
    Since \(X \geq 0\) and \(X \in \NN\), we have
    \[
        X = \sum_{k = 1}^{\infty} \indicator \para{X \geq k} = \sum_{k = 0}^{\infty} \indicator \para{X > k}.
    \]

    Indeed, this is true for any \(X = x \in \NN\), since
    \[
        x = \sum_{k = 1}^{x} 1 = \sum_{k = 1}^{x} \indicator \para{x \geq k} = \sum_{k = 1}^{\infty} \indicator \para{x \geq k}
    \]
    and similarly for the strict inequality case.

    Hence,
    \[
        X \para{\omega} = \sum_{k = 1}^{\infty} \indicator \para{X \para{\omega} \geq k} = \sum_{k = 0}^{\infty} \indicator \para{X \para{\omega} > k}.
    \]

    Taking the expectation, we get
    \begin{align*}
        \Expt\brac{X} &= \Expt\brac{\sum_{k = 1}^{\infty} \indicator\para{X \geq k}} \\
        &= \sum_{k = 1}^{\infty} \Expt\brac{\indicator\para{X \geq k}} \\
        &= \sum_{k = 1}^{\infty} \Prob\para{X \geq k},
    \end{align*}
    and similarly for the strict case.
\end{proof}

The indicator random variable gives an alternative proof of \autoref{thm:inclusion-exclusion} (Inclusion--Exclusion Formula).

\begin{proof}[of \autoref{thm:inclusion-exclusion} (Inclusion--Exclusion Formula), using Indicator Functions]
    First, note that for events \(A, B \in \mathcal{F}\),
    \begin{itemize}
        \item \(\indicator\para{A\compl} = 1 - \indicator\para{A}\);
        \item \(\indicator\para{A \cap B} = \indicator\para{A} \cdot \indicator\para{B}\); and for the union,
        \item \(\indicator\para{A \cup B} = 1 - \para{1 - \indicator\para{A}} \para{1 - \indicator\para{B}}\).
    \end{itemize}

    More generally, for any \(A_1, A_2, \ldots, A_n \in \mathcal{F}\), we have
    \begin{align*}
        \indicator\para{A_1 \cup A_2 \cup \cdots \cup A_n} &= 1 - \prod_{i = 1}^{n} \para{1 - \indicator\para{A_i}} \\
        &= \sum_{i = 1}^{n} \indicator \para{A_i} - \sum_{i_1 < i_2} \indicator\para{A_{i_1} \cap A_{i_2}} + \cdots + \para{-1}^{n + 1} \indicator\para{A_1 \cap \cdots \cap A_n}.
    \end{align*}

    We may take the expectation on both sides, and use that
    \[
        \Expt\brac{\indicator\para{A}} = \Prob\para{A}.
    \]

    Hence,
    \[
        \Prob\para{A_1 \cup A_2 \cup \cdots \cup A_n} = \sum_{k = 1}^{n} \sum_{1 \leq i_1 < i_2 < \cdots < i_k \leq n} \Prob\para{A_{i_1} \cap A_{i_2} \cap \cdots \cap A_{i_n}}.
    \]
\end{proof}

\subsection{Variance}

Defining the expectation allows us to define several other important things about a random variable.

\begin{definition}[Moment]
    \label{def:moment}%
    Let \(X\) be a random variable, and \(r \in \NN\). \(\Expt\brac{X^r}\), i well-defined, is the \emph{\(r\)-th moment} of \(X\).
\end{definition}

\begin{definition}[Variance]
    \label{def:variance}%
    The \emph{variance} of \(X\) is given as
    \[
        \Var\para{X} = \Expt\brac{\para{X - \Expt\brac{X}}^2}.
    \]
\end{definition}

The variance is a `measure' of how well concentrated \(X\) is around its expectation. The smaller the variance, the more concentrated \(X\) is.

\begin{definition}[Standard Deviation]
    \label{def:std-deviation}%
    \(\sqrt{\Var\para{X}}\) is known as the \emph{standard deviation}.
\end{definition}

\begin{proposition}[Properties of the Variance]
    \label{prop:properties-of-variance}%
    \begin{itemize}
        \item \(\Var\para{X} \geq 0\).
        \item If \(\Var\para{X} = 0\), then \(\Prob\para{X = \Expt\brac{X}} = 1\).
        \item If \(c \in \RR\), then \(\Var\para{cX} = c^2 \Var\para{X}\), and \(\Var\para{X + c} = \Var\para{X}\).
    \end{itemize}
\end{proposition}

\begin{proposition}[Alternative Formulation of Variance]
    \label{prop:alt-formulation}%
    We have
    \[
        \Var\para{X} = \Expt\brac{X^2} - \Expt\brac{X}^2.
    \]
\end{proposition}

\begin{proof}
    We have
    \begin{align*}
        \Var\para{X} &= \Expt\brac{X^2 - 2X \Expt\brac{X} + \para{\Expt\brac{X}}^2} \\
        &= \Expt\brac{X^2} - 2 \Expt\brac{X} \Expt\brac{X} + \para{\Expt\brac{X}}^2 \\
        &= \Expt\brac{X^2} - \para{\Expt\brac{X}}^2.
    \end{align*}
\end{proof}

\begin{proposition}[Variance as a Minimalisation]
    \label{prop:variance-minimalisation}%
    We have
    \[
        \Var\para{X} = \min_{x \in \RR} \Expt\brac{\para{X - c}^2},
    \]
    obtained for \(c = \Expt\brac{X}\).
\end{proposition}

\begin{proof}
    We let
    \[
        f\para{c} = \Expt\brac{\para{X - c}^2}.
    \]

    Then, we have
    \[
        f\para{c} = \Expt\brac{X^2} = 2c \Expt\brac{X} + c^2.
    \]

    Hence, differentiating gives that
    \[
        f'\para{c} = -2 \Expt\brac{X} + 2c
    \]
    and hence
    \[
        \min_{c \in \RR} f\para{c} = f\para{\Expt\para{X}} = \Var\para{X}.
    \]
\end{proof}

\begin{examples}[Variances of Common Random Variables]
    \label{ex:variance-rv}%
    \begin{enumerate}
        \item If \(X \sim \Binomial\para{n, p}\), then \(\Expt\brac{X} = np\).
        
        We have
        \[
            \Var\para{X} = \Expt\brac{X\para{X - 1}} + \Expt\brac{X} - \para{\Expt\brac{X}}^2.
        \]

        We can find
        \begin{align*}
            \Expt\brac{X \para{X - 1}} &= \sum_{k = 2}^{n} k\para{k - 1} \cdot \frac{n!}{k! \para{n - k}!} \cdot p^k \cdot \para{1 - p}^{n - k} \\
            &= \para{n - 1} n p^2 \sum_{k = 0}^{n - 2} \binom{n - 2}{k} p^k \para{1 - p}^{n - 2 - k} \\
            &= \para{n - 1} n p^2.
        \end{align*}

        Substituting this back, we have
        \[
            \Var\para{X} = \para{n - 1} np^2 + np - \para{np}^2 = np \para{1 - p}.
        \]

        (We will see that using that a binomial random variable is the sum of independent Bernoulli random variables, we may find this using that variance is linear for independent variables.)

        \item Let \(X \sim \Poisson\para{\lambda}\) for \(\lambda > 0\), and recall that \(\Expt\brac{X} = \lambda\).
        
        We have
        \begin{align*}
            \Expt\brac{X \para{X - 1}} &= \sum_{k = 2}^{\infty} k \para{k - 1} e^{-\lambda} \frac{\lambda^k}{k!} \\
            &= e^{-\lambda} \lambda^2 \sum_{k = 2}^{\infty} \frac{\lambda^{k - 2}}{\para{k - 2}!} \\
            &= \lambda^2
        \end{align*}
        and hence
        \[
            \Var\para{X} = \lambda^2 + \lambda - \lambda^2 = \lambda.
        \]
    \end{enumerate}
\end{examples}

\subsection{Covariance}

\begin{definition}[Covariance]
    \label{def:covariance}%
    Let \(X\) and \(Y\) be random variables. We let
    \[
        \Cov\para{X, Y} = \Expt\brac{\para{X - \Expt\brac{X}} \para{Y - \Expt\brac{Y}}}
    \]
    be the \emph{covariance} of \(X\) and \(Y\).
\end{definition}

It is a `measure' of how dependent \(X\) and \(Y\) are.

\begin{proposition}[Properties of Covariance]
    \label{prop:properties-of-covariance}%
    \begin{itemize}
        \item \(\Cov\para{X, Y} = \Cov\para{Y, X}\). 
        \item \(\Cov\para{X, X} = \Var\para{X}\).
        \item \(\Cov\para{X, Y} = \Expt\brac{X Y} - \Expt\brac{X} \Expt\brac{Y}\).
        \item For \(c in \RR\), \(\Cov\para{cX, Y} = c \Cov\para{X, Y}\), and \(\Cov\para{c + X, Y} = \Cov\para{X, Y}\).
        \item For \(c \in \RR\), \(\Cov\para{c, X} = 0\).
    \end{itemize}
\end{proposition}

This gives us a nice formula on the variance of a sum as well.

\begin{proposition}[Variance of Sum]
    \label{prop:variance-of-sum}%
    We have
    \[
        \Var\para{X + Y} = \Var\para{X} + \Var\para{Y} + 2 \Cov\para{X, Y}.
    \]    
\end{proposition}

\begin{proof}
    We have
    \begin{align*}
        \Var\para{X + Y} &= \Expt\brac{\para{\para{X - \Expt\para{X}} + \para{Y - \Expt\para{Y}}}^2} \\
        &= \Var\para{X} + 2 \Cov\para{X, Y} + \Var\para{Y}.
    \end{align*}
\end{proof}

\begin{proposition}[Linearity of Covariance]
    \label{prop:linearity-of-covariance}%
    For random variables \(X, Y\) and \(Z\), we have
    \[
        \Cov\para{X + Y, Z} = \Cov\para{X, Z} + \Cov\para{Y, Z}.
    \]

    More generally, we for any real constants  \(c_1, c_2, \ldots, c_n\) and \(d_1, d_2, \ldots, d_m\) and random variables \(X_1, X_2, \ldots, X_n\) and \(Y_1, Y_2, \ldots, Y_m\), that
    \[
        \Cov\para{\sum_{i = 1}^{n} c_i X_i, \sum_{j = 1}^{m} d_j Y_j} = \sum_{i = 1}^{n} \sum_{j = 1}^{m} c_i d_j \Cov\para{X_i, Y_j}.
    \]
\end{proposition}

\begin{corollary}[Variance of Sum]
    \label{cor:variance-of-sum}%
    We have
    \[
        \Var\para{\sum_{i = 1}^{n} X_i} = \Cov\para{\sum_{i = 1}^{n} X_i, \sum_{i = 1}^{n} X_i} = \sum_{i = 1}^{n} \Var\para{X_i} + \sum_{i \neq j} \Cov\para{X_i, X_j}.
    \]
\end{corollary}

We stated before that the covariance is a `measure' of how independent two random variables are. We first see some important properties on independence that will help with this.

\begin{lemma}[Independence implies Pairwise Independence]
    \label{lem:indep-pairwise-indep}%
    Let \(X_1, X_2, X_3\) be independent, i.e.,
    \[
        \Prob\para{X_1 = x_1, X_2 = x_2, X_3 = x_3} = \prod_{i = 1}^{3} \Prob\para{X_i = x_i}.
    \]

    Then \(X_1\) is independent of \(X_2\).
\end{lemma}

\begin{proof}
    We have
    \begin{align*}
        \Prob\para{X_1 = x_1, X_2 = x_2} &= \sum_{x_3} \Prob\para{X_1 = x_1, X_2 = x_2, X_3 = x_3} \\
        &= \sum_{x_3} \Prob\para{X_1 = x_1} \Prob\para{X_2 = x_2} \Prob\para{X_3 = x_3} \\
        &= \Prob\para{X_1 = x_1} \Prob\para{X_2 = x_2} \sum_{x_3} \Prob\para{X_3 = x_3} \\
        &= \Prob\para{X_1 = x_1} \Prob\para{X_2 = x_2}.
    \end{align*}
\end{proof}

\begin{lemma}[Independence and Expectation of Product]
    \label{lem:independence-expectation-product}%
    Let \(X\) and \(Y\) be independent random variables, and \(\functiondomain{f}{g}{\RR}{\RR_{+}}\). Then
    \[
        \Expt\brac{f\para{X} g\para{Y}} = \Expt\brac{f\para{X}} \Expt\brac{g\para{Y}}.
    \]
\end{lemma}

\begin{proof}
    Let \(Z = \para{X, Y}\), and define
    \[  
        h\para{Z} = f\para{X} g\para{Y}.
    \]

    Then we have
    \begin{align*}
        \Expt\brac{h\para{Z}} &= \sum_{\para{x, y}} h\para{x, y} \Prob\para{Z = \para{x, y}} \\
        &= \sum_{\para{x, y}} f\para{x} g\para{y} \Prob\para{X = x, Y = y} \\
        &= \sum_{x, y} f\para{x} g\para{y} \Prob\para{X = x} \Prob\para{Y = y} \\
        &= \sum_{x} f\para{x} \Prob\para{X = x} \sum_{y} g\para{y} \Prob\para{Y = y} \\
        &= \Expt\brac{f\para{X}} \cdot \Expt\brac{g\para{Y}}.
    \end{align*}
\end{proof}

\begin{proposition}[Independence gives Zero Covariance]
    If \(X\) and \(Y\) are independent, then
    \[
        \Cov\para{X, Y} = 0.
    \]
\end{proposition}

\begin{proof}
    This follows from \autoref{lem:independence-expectation-product}, since
    \[
        \Expt\brac{X - \Expt\brac{X}} = \Expt\brac{Y - \Expt\brac{Y}} = 0.
    \]
\end{proof}

The converse does not hold -- in general, for two random variables \(X\) and \(Y\), \(\Cov\para{X, Y} = 0\) does \textbf{not} tell us that \(X\) and \(Y\) are independent.

\begin{example}[Zero Variance but Not Independent]
    \label{ex:zero-variance-no-independent}%
    Let \(X_1, X_2, X_3 \sim \Bernoulli\para{\frac{1}{2}}\). Set
    \[
        Y_1 = 2 X_1 - 1, \quad Y_2 = 2 X_2 - 1,
    \]
    and
    \[
        Z_1 = Y_1 \cdot X_3, \quad Z_2 = Y_2 \cdot X_3.
    \]

    Then
    \[
        \Expt\brac{Y_1} = \Expt\brac{Y_2} = 0.
    \]

    By independence,
    \[
        \Expt\brac{Z_1} = \Expt\brac{Y_1 \cdot X_3} = \Expt\brac{Y_1} \cdot \Expt\brac{X_3} = 0,
    \]
    and similarly \(\Expt\brac{Z_2} = 0\).

    Therefore,
    \[
        \Cov\para{Z_1, Z_2} = \Expt\brac{Z_1 \cdot Z_2} = \Expt\brac{Y_1 \cdot Y_2 \cdot X_3^2} = 0.
    \]

    We want to show \(Z_1\) is not independent of \(Z_2\). Indeed,
    \[
        \Prob\para{Z_1 = 0, Z_2 = 0} = \Prob{X_3 = 0} = \frac{1}{2},
    \]
    but
    \[
        \Prob\para{Z_1 = 0} = \Prob\para{Z_2 = 0} = \frac{1}{2}.
    \]
\end{example}